\documentclass[a4paper,11pt]{article}
\usepackage[utf8]{inputenc}
\usepackage[T1]{fontenc}
\usepackage[french]{babel}
\usepackage{geometry}
\usepackage{enumitem}
\usepackage{titlesec}
\usepackage{xcolor}
\usepackage{hyperref}
\usepackage{booktabs}
\usepackage{array}
\usepackage{amsmath}

\geometry{margin=2.5cm}

\definecolor{primary}{RGB}{91, 92, 246}
\definecolor{done}{RGB}{34, 197, 94}
\definecolor{partial}{RGB}{245, 158, 11}

\titleformat{\section}{\large\bfseries\color{primary}}{}{0em}{}[\titlerule]
\titleformat{\subsection}{\normalsize\bfseries}{}{0em}{}

\hypersetup{colorlinks=true, linkcolor=primary, urlcolor=primary}

\title{\textbf{TravelPath — Fonctionnalités implémentées}\\[4pt]
\large Application Android · Java · Firebase · minSdk\,24}
\author{}
\date{}

\begin{document}
\maketitle
\tableofcontents
\newpage

% ──────────────────────────────────────────────────────────────────
\section{Saisie des préférences}

\subsection{Activités}
L'utilisateur sélectionne une ou plusieurs activités via des chips multi-sélection.
Chaque activité est traduite en catégories OpenTripMap (\texttt{kinds}) pour orienter
la recherche de points d'intérêt.

\begin{center}
\begin{tabular}{ll}
  \toprule
  Activité & Catégories OpenTripMap associées \\
  \midrule
  Culture      & \texttt{museums}, \texttt{theatres\_and\_entertainments},
                 \texttt{historic}, \texttt{architecture} \\
  Restauration & \texttt{foods}, \texttt{restaurants}, \texttt{cafes} \\
  Nature       & \texttt{natural}, \texttt{gardens\_and\_parks}, \texttt{beaches} \\
  Loisirs      & \texttt{amusements}, \texttt{sport}, \texttt{shops} \\
  Découverte   & \texttt{tourist\_facilities}, \texttt{monuments\_and\_memorials},
                 \texttt{view\_points}, \texttt{historic} \\
  \bottomrule
\end{tabular}
\end{center}

\subsection{Liste de lieux favoris / étapes obligatoires}
L'utilisateur peut ajouter manuellement des étapes via une boîte de dialogue
(nom, description, photo depuis la galerie) avant ou après la génération automatique.
Ces étapes s'intègrent au parcours final aux côtés des étapes générées.

\subsection{Budget maximum}
Un curseur permet de définir un budget total en euros.
La contrainte est appliquée de manière cumulative lors de la sélection des étapes :
chaque lieu candidat est écarté si son coût d'entrée estimé dépasse le budget restant.
Les coûts estimés sont : restaurants 18\,€, musées 12\,€, cinémas/spectacles 10\,€,
attractions 8\,€, cafés 5\,€, monuments/parcs/points de vue 0\,€.

\subsection{Durée}
Un curseur de 1 à 8 heures détermine trois paramètres dérivés :
\begin{itemize}[leftmargin=1.5em]
  \item le budget-temps total (durée $\times$ 60 minutes) utilisé comme
        contrainte dure lors de la sélection des étapes ;
  \item le nombre maximum d'étapes candidates
        $\text{maxSteps} = \lfloor\text{durée}_{h} \times 60 / 20\rfloor$ ;
  \item le rayon de recherche des POI
        $r = \max(500,\; \text{maxMarcheKm}/2 \times 1000)$\,m.
\end{itemize}

\subsection{Effort / tolérances météo}

\paragraph{Niveau d'effort.}
Trois niveaux sont disponibles (\textit{Facile}, \textit{Modéré},
\textit{Difficile}), modulant la fraction du temps consacrée à la marche
(respectivement 30\,\%, 40\,\%, 60\,\%) et donc le rayon de recherche.

\paragraph{Sensibilité météo.}
Deux options météo sont disponibles :
\begin{itemize}[leftmargin=1.5em]
  \item \textit{Sensible à la pluie} : seuls les lieux couverts (musées,
        théâtres, restaurants, cafés, commerces) sont retenus quand
        le code météo indique des précipitations ($\geq 51$).
  \item \textit{Sensible à la chaleur} : les lieux extérieurs exposés
        (plages, points de vue, jardins) sont exclus quand
        la température dépasse 30\,°C.
\end{itemize}
La météo en temps réel est fournie par l'\textbf{API Open-Meteo}
(\texttt{current\_weather}) pour la ville saisie.

% ──────────────────────────────────────────────────────────────────
\section{Génération de parcours — 3 options}

\subsection{Algorithme de sélection des étapes}

La génération repose sur un algorithme glouton (\textit{chain-based greedy})
qui construit le parcours étape par étape, depuis la dernière étape ajoutée.

\paragraph{Score de chaque candidat.}
\[
  \text{score} =
    \underbrace{\text{rate} \times 20}_{\text{popularité OTM}}
  + \underbrace{\text{activityMatch} \times 12}_{\text{pertinence activités}}
  + \underbrace{\text{spreadScore}}_{\text{étalement géographique}}
\]
avec
\[
  \text{spreadScore} =
  \begin{cases}
    0    & \text{si première étape} \\
    -20  & \text{si } d < 0{,}3\,\text{km} \quad (\text{micro-cluster pénalisé}) \\
    20 - \lvert d - d_{\text{cible}} \rvert \times 4
         & \text{si } d \leq 1{,}5\,d_{\text{cible}} \quad (\text{plage idéale}) \\
    5 - (d - 1{,}5\,d_{\text{cible}}) \times 6
         & \text{sinon} \quad (\text{détour pénalisé})
  \end{cases}
\]
où $d_{\text{cible}} = \text{maxRouteKm} / \text{maxSteps}$ est la distance
idéale par tronçon pour répartir les étapes sur toute l'aire disponible.

\paragraph{Contraintes dures.}
\begin{itemize}[leftmargin=1.5em]
  \item Distance minimale de 150\,m entre toute paire d'étapes (anti-clustering).
  \item Budget cumulatif : $\text{coût} + \text{frais\_étape} \leq \text{budget\_restant}$.
  \item Temps cumulatif :
        $\lceil d_{\text{dernière}} \times 12 \rceil + \text{duréeVisite} \leq \text{tempsRestant}$
        (marche à 5\,km/h = 12\,min/km).
  \item Distance totale : écrêtage si la route dépasse \texttt{maxRouteKm}.
\end{itemize}

\paragraph{Durée de visite estimée par type de lieu.}
\begin{center}
\begin{tabular}{ll}
  \toprule
  Type de lieu & Durée estimée \\
  \midrule
  Théâtres / Cinémas           & 120 min \\
  Musées                       & 90 min \\
  Plages                       & 90 min \\
  Restaurants / Gastronomie    & 60 min \\
  Attractions / Loisirs        & 60 min \\
  Parcs / Jardins              & 40 min \\
  Lieux religieux              & 25 min \\
  Points de vue                & 20 min \\
  Monuments / Architecture     & 15--20 min \\
  Cafés                        & 30 min \\
  Défaut                       & 20 min \\
  \bottomrule
\end{tabular}
\end{center}

\paragraph{Optimisation de l'ordre des visites (TSP).}
\begin{itemize}[leftmargin=1.5em]
  \item \textit{Branch-and-bound} exact pour $N \leq 9$ étapes.
  \item \textit{Nearest-Neighbor + 2-opt} pour $N > 9$.
\end{itemize}

\subsection{Trois profils proposés}

\begin{center}
\begin{tabular}{lccc}
  \toprule
  Profil & Filtrage des lieux & Frais fixes / étape & Budget \\
  \midrule
  \textbf{Économique} & Entrée $< 10$\,€ uniquement & 3\,€ & 100\,\% \\
  \textbf{Équilibré}  & Tous les lieux              & 5\,€ & 100\,\% \\
  \textbf{Confort}    & Tous les lieux (premium)    & 10\,€ & 130\,\% \\
  \bottomrule
\end{tabular}
\end{center}

Chaque option est présentée sur une carte avec couleur thématique, icône,
badge \textit{Recommandé} sur l'option centrale, nombre d'étapes,
distance totale + durée estimée, budget estimé et séquence des noms d'étapes.

% ──────────────────────────────────────────────────────────────────
\section{Itinéraire et étapes}

\subsection{Ordre des visites}
L'ordre optimal des étapes est calculé par l'algorithme TSP décrit
à la section précédente, minimisant la distance totale de marche.

\subsection{Distances et temps}
La durée affichée est la somme réelle des durées de visite par type de lieu
et des temps de marche entre étapes consécutives (en minutes), convertie en
format « 4h30 » ou « 45 min ».
La distance totale (en km) est calculée sur le chemin TSP-optimisé.

\subsection{Suggestions de créneaux horaires (ouverture)}
Les horaires d'ouverture réels sont récupérés via l'\textbf{API Overpass}
(tag OSM \texttt{opening\_hours}, rayon 100\,m autour de chaque étape) et
affichés dans la fiche détaillée de l'étape.

% ──────────────────────────────────────────────────────────────────
\section{Présentation riche}

\subsection{Carte interactive}
\begin{itemize}[leftmargin=1.5em]
  \item Carte OSMDroid (tuiles Mapnik) avec marqueurs numérotés sur chaque étape.
  \item Tracé de l'itinéraire pédestre via \textbf{OpenRouteService}
        (profil \texttt{foot-walking}, requête POST GeoJSON) ; fallback ligne
        droite en cas d'indisponibilité de l'API.
  \item Carte miniature intégrée dans la page de présentation du parcours.
\end{itemize}

\subsection{Galerie photos par étape}
\begin{itemize}[leftmargin=1.5em]
  \item Photos récupérées automatiquement via \textbf{OpenTripMap XID}
        (\texttt{preview.source}) et \textbf{Wikipedia Search API}
        (fr puis en, miniature 600\,px).
  \item Photo manuelle possible à la création (galerie ou appareil photo).
  \item Affichage dans la liste des étapes et dans la fiche bottom sheet.
  \item Chargement parallèle via un pool de 4 threads ; URL stockée dans
        Firestore (pas de base64 pour rester sous la limite de 1\,Mo/document).
\end{itemize}

\subsection{Fiche détaillée par étape (bottom sheet)}
Un appui sur une étape ouvre un panneau glissant contenant :
\begin{itemize}[leftmargin=1.5em]
  \item En-tête image (220\,dp) avec dégradé, badge numéro et poignée de glissement.
  \item Prix d'entrée (tag OSM \texttt{fee}/\texttt{charge} via Overpass, rayon 150\,m).
  \item Horaires d'ouverture (tag \texttt{opening\_hours} via Overpass).
  \item Description (extrait Wikipédia ou champ \texttt{info.descr} d'OpenTripMap).
  \item Bouton \textit{Voir sur la carte} ouvrant l'intent \texttt{geo:lat,lon}.
\end{itemize}

\subsection{Résumé et métriques du parcours}
La page de présentation affiche :
\begin{itemize}[leftmargin=1.5em]
  \item Titre, ville, description, auteur.
  \item Chips métriques : durée, budget, difficulté.
  \item Résumé météo temps réel (température + code météo iconique).
  \item Photo de couverture (base64 Firestore ou image ressource).
\end{itemize}

% ──────────────────────────────────────────────────────────────────
\section{Interaction}

\subsection{Like / délike d'un parcours}
\begin{itemize}[leftmargin=1.5em]
  \item Bouton like avec compteur mis à jour instantanément dans l'UI
        et persisté dans Firestore (\texttt{likeCount}).
  \item État \textit{liked} mémorisé par utilisateur ; toggle cohérent
        (like $\to$ délike décrémente le compteur).
\end{itemize}

\subsection{Mise en favoris}
\begin{itemize}[leftmargin=1.5em]
  \item Ajout/retrait d'un parcours ou d'une photo en favoris,
        avec compteur \texttt{favoriteCount} dans Firestore.
  \item Fragment \textit{Favoris} dédié avec onglets Parcours / Photos.
\end{itemize}

\subsection{Regénération avec ajustements}
L'utilisateur peut modifier n'importe quel paramètre (durée, budget, activités,
effort, météo) et relancer la génération pour obtenir de nouvelles options,
sans limite de nombre de tentatives.

\subsection{Sauvegarde}
Le parcours sélectionné est sauvegardé dans \textbf{Cloud Firestore}
avec toutes ses étapes (nom, coordonnées, description, URL d'image, XID).

\subsection{Partage}
\begin{itemize}[leftmargin=1.5em]
  \item Partage du fichier PDF généré via l'intent Android \texttt{ACTION\_SEND}
        (\texttt{FileProvider}) vers n'importe quelle application.
  \item Visibilité public/privé configurable à la publication.
\end{itemize}

\subsection{Publication planifiée}
\begin{itemize}[leftmargin=1.5em]
  \item Sélecteur date + heure pour programmer une publication future.
  \item Exécution par \textbf{WorkManager} (\texttt{ScheduledPublishWorker})
        avec persistance entre redémarrages.
  \item Notification automatique des abonnés aux tags d'activités lors
        de la publication (push Firestore).
\end{itemize}

% ──────────────────────────────────────────────────────────────────
\section{Mode hors-ligne léger}

\begin{itemize}[leftmargin=1.5em]
  \item \textbf{Cache tuiles OSMDroid} : les tuiles de la carte consultées
        sont stockées sur le stockage interne du device
        (\texttt{CacheDir/osmdroid/tiles}), permettant de revoir la carte
        sans connexion.
  \item \textbf{Cache mémoire des parcours} : les parcours chargés depuis
        Firestore sont conservés dans des registres en mémoire
        (\texttt{PathRegistry}, \texttt{ActivePathRegistry}, \texttt{PathRepository})
        pour la durée de la session, évitant les requêtes répétées.
  \item \textbf{Images en base64} : les photos de couverture et photos manuelles
        des étapes sont encodées en base64 et stockées directement dans
        Firestore, accessibles sans requête réseau supplémentaire.
\end{itemize}

% ──────────────────────────────────────────────────────────────────
\section{Prise en compte de la météo et des créneaux d'ouverture}

\subsection{Météo du jour}
\begin{itemize}[leftmargin=1.5em]
  \item Appel à \textbf{Open-Meteo} (\texttt{/v1/forecast?current\_weather=true})
        pour la ville géocodée, au moment de la génération.
  \item Récupération de la température (°C) et du code météo WMO.
  \item Affichage d'un résumé iconique : \textit{☀ 22°C}, \textit{🌧 14°C}, etc.
  \item Filtrage des POI inadaptés selon les préférences météo de l'utilisateur
        (cf. section\,1.5).
\end{itemize}

\subsection{Créneaux d'ouverture}
\begin{itemize}[leftmargin=1.5em]
  \item Interrogation de l'\textbf{API Overpass} (tag \texttt{opening\_hours})
        au rayon de 100\,m autour de chaque étape, à l'affichage de la
        fiche détaillée.
  \item Valeur brute OSM affichée telle quelle (ex. \texttt{Mo-Fr 09:00-18:00}).
  \item Interrogation du tag \texttt{fee} / \texttt{charge} pour afficher
        le prix d'entrée réel.
\end{itemize}

% ──────────────────────────────────────────────────────────────────
\section{Export du parcours en PDF}

La génération du PDF est réalisée nativement avec \texttt{android.graphics.pdf.PdfDocument}
(pas de dépendance externe), dans un thread d'arrière-plan.

\subsection{Contenu du document}
\begin{itemize}[leftmargin=1.5em]
  \item \textbf{En-tête} : nom de l'application, titre du parcours, ville, auteur,
        description.
  \item \textbf{Chips métriques} : durée, budget, difficulté — rendu sous forme
        de badges colorés sur fond \texttt{\#EEF2FF}.
  \item \textbf{Fiche par étape} : numéro, nom, description, photo (Bitmap)
        redimensionnée en 120\,$\times$\,90\,px.
  \item Saut de page automatique quand la hauteur restante est insuffisante
        pour la prochaine étape.
\end{itemize}

\subsection{Partage}
Le PDF est écrit dans le stockage cache de l'application, puis partagé
via un \texttt{FileProvider} et l'intent \texttt{ACTION\_SEND}
(\texttt{application/pdf}).

% ──────────────────────────────────────────────────────────────────
\section{Enrichissement automatique des étapes}

\subsection{À la création}
\begin{enumerate}[leftmargin=1.5em]
  \item \textbf{OpenTripMap XID} : description (\texttt{wikipedia\_extracts}
        ou \texttt{info.descr}) et URL d'image (\texttt{preview.source}
        ou champ \texttt{image}).
  \item \textbf{Wikipedia REST} (fr puis en,
        \texttt{/api/rest\_v1/page/summary/\{titre\}}) : extrait textuel
        et miniature — fallback si OTM n'a pas de données.
  \item \textbf{Overpass} (rayon 150\,m, tag \texttt{fee}) : prix d'entrée
        réel depuis les données OSM.
\end{enumerate}

\subsection{À l'ouverture d'un parcours existant}
Pour les étapes sans image (parcours créés avant la fonctionnalité) :
\begin{enumerate}[leftmargin=1.5em]
  \item Recherche OTM par coordonnées (rayon 80\,m).
  \item Fallback Wikipedia Search API
        (\texttt{action=query\&generator=search}) — plus robuste que
        l'API summary pour les titres approximatifs.
\end{enumerate}
Chargement parallèle (pool de 4 threads) ; l'URL est mémorisée en mémoire
pour la session.

\end{document}
